\begin{textsample}{POS Dim 4 – 2014\_c5 – Score 1.00 – 2014\_c5/enade\_2014\_eng\_civ\_231731472f.txt}  \label{ex:f4_pos_159}
CAPÍTULO 5 DISTRIBUIÇÃO DOS CONCEITOS 5.1 PANORAMA NACIONAL DA DISTRIBUIÇÃO DOS CONCEITOS O Conceito Enade é um indicador de qualidade que avalia o desempenho dos estudantes a partir dos resultados do ENADE.
Os valores possíveis vão de 1 ( pior situação ) a 5 ( melhor situação ).
As Unidades de \textbf{Observação} com apenas um ou sem nenhum concluinte participante não obtêm o Conceito Enade, ficando Sem Conceito ( SC ).
A Tabela 5.1 apresenta a quantidade e distribuição de cursos de Engenharia Civil participantes do ENADE/2014, por faixa de conceito e Grande Região.
Mais uma vez, ressaltando que a diferença entre os cursos tabulados neste capítulo e no capítulo 2 são os cursos sem conceito, em princípio, aqueles sem alunos concluintes que participassem da prova ou que tivessem tão somente um aluno concluinte21.
Observando -se os dados da Tabela 5.1, nota -se que, dos 283 cursos participantes, 103 ( 36,4 \% ) classificaram -se com conceito 2, o valor modal.
Este foi também o conceito modal em mais três regiões : Norte ( 35,3 \% ), Nordeste ( 36,2 \% ) e Sudeste ( 40,1 \% ).
Nas regiões Sul e Centro-Oeste, a moda foi o conceito 3, atribuído a 36,4 \% dos cursos situados em cada uma dessas regiões.
O conceito 3 foi o segundo mais frequente em nível nacional ( 31,1 \%, correspondendo a 88 cursos ) e o conceito 4, o terceiro ( 19,4 \%, correspondendo a 55 cursos ).
Houve, ainda, 12 cursos ( 4,2 \% ) que receberam conceito 5 e 23 cursos ( 8,1 \% ) que receberam conceito 1.
De todos os 283 cursos de Engenharia Civil, apenas dois cursos ( 0,7 \% ) ficaram sem conceito ( SC ). 21 Estes cursos com somente um concluinte não foram considerados no Capítulo 2.
A região Norte participou com 17 cursos ou 6,0 \% do total nacional.
Nesta região, o conceito 2 foi atribuído a seis cursos dos 17 participantes, o que equivale a 35,3 \% do total regional, além de ser o conceito modal na região.
Os conceitos 1 e 3 foram atribuídos a quatro cursos ( 23,5 \% ), cada.
Dois cursos receberam o conceito 4 e um curso, o conceito 5.
Nenhum curso dessa região ficou sem conceito ( SC ).
A região Nordeste participou com 47 cursos ou 16,6 \% do total nacional.
Nesta região, 17 cursos, 36,2 \% em termos regionais, obtiveram conceito 2, o conceito modal para a região.
Os conceitos 1 e 3 foram atribuídos a, respectivamente, dois e 11 cursos ( 4,3 \% e 23,4 \% ). 15 cursos ( 31,9 \% ) obtiveram conceito 4 e dois cursos ( 2,8 \% ) receberam conceito 5.
Nessa região, nenhum curso ficou sem conceito ( SC ).
Dos 142 cursos participantes da região Sudeste ( 50,2 \% do total nacional ), o conceito 2 também foi o conceito modal, tendo sido o obtido por 57 ( 40,1 \% ) cursos.
O conceito 1 foi atribuído a 14 cursos ( 9,9 \% ) e o conceito 3, a 45 cursos ( 31,7 \% ).
Os conceitos 4 e 5 foram atribuídos, respectivamente, a 20 cursos ( 14,1 \% ) e a quatro cursos ( 2,8 \% ).
Os dois cursos de Engenharia Civil que ficaram sem conceito ( SC ) pertencem a região Sudeste ( 1,4 \% dos cursos da região ).
Nos 55 cursos da região Sul, como já comentado, há predominância do conceito 3, atribuído a 20 cursos, 36,4 \% dos cursos da região ( conceito modal ).
O conceito 1 foi atribuído a um curso ( 1,8 \% ) e o conceito 2, a 17 cursos ( 30,9 \% ).
O conceito 4 foi recebido por 14 cursos ( 25,5 \% ) e o conceito 5, por três cursos ( 5,5 \% ).
Nenhum curso na região Sul ficou sem conceito ( SC ).
Já dos 22 cursos participantes na região Centro-Oeste, oito receberam conceito 3, o conceito modal ( 36,4 \% ).
Nesta região, nenhum curso ficou sem conceito ( SC ).
Em cada um dos conceitos 1 e 5 foram alocados dois cursos ( 9,1 \% cada ).
Seis cursos ( 27,3 \% ) receberam conceito 2, correspondendo a 27,3 \% dos cursos nessa região, e quatro cursos ( 18,2 \% ) foram avaliados com conceito 4.
O Gráfico 5.1 apresenta a distribuição do Conceito Enade segundo Grande Região : poligonais mais à esquerda representam Grandes Regiões com pior distribuição, e poligonais mais à direita, Grandes Regiões com melhores conceitos.
As regiões Norte ( linha azul ) e Sudeste ( linha vermelha ) apresentam os cursos em IES com conceitos mais baixos, ao passo que as regiões e Sul ( linha roxa ), Nordeste ( linha verde ) e Centro-Oeste ( linha preta ) apresentam as distribuições com valores maiores. 5.2 CONCEITOS POR CATEGORIA ADMINISTRATIVA E POR GRANDE REGIÃO A Tabela 5.2 apresenta a distribuição dos cursos participantes do ENADE/2014 de Engenharia Civil, por Categoria Administrativa, de acordo com os conceitos por eles alcançados, segundo as Grandes Regiões brasileiras.
Dos 283 cursos participantes, 91 ( 32,2 \% ) eram ministrados em Instituições Públicas e 192 ( 67,8 \% ), em Instituições Privadas.
De acordo com as informações da Tabela 5.2, em termos nacionais, entre Instituições Públicas e Privadas : das 12 instituições que receberam o conceito 5, 11 eram Públicas e apenas uma, Privada.
Dos 91 cursos participantes de IES Públicas, o conceito 4 foi o valor modal, atribuído a 44 cursos ( 48,4 \% ).
Entre os demais cursos de Instituições Públicas participantes, um curso ( 1,1 \% da categoria ) foi avaliado com conceito 1, sete cursos ( 7,7 \% ) obtiveram conceito 2, 28 cursos ( 30,8 \% ) receberam conceito 3 e, como já comentado, 11 cursos foram alocados ao conceito 5 ( 12,1 \% ).
Nesta categoria, nenhum curso ficou sem conceito ( SC ).
Na rede privada, o conceito modal foi 2, atribuído a 96 ( 50,0 \% ) dos 192 cursos da categoria.
Entre os demais cursos participantes, 22 receberam conceito 1 e 60, conceito 3.
O conceito 4 foi atribuído a 11 cursos e o conceito 5, a um único curso.
Nesta Categoria Administrativa, dois cursos ficaram sem conceito ( SC ).
Na análise por região, observa -se que, na região Norte, as Instituições Privadas participaram com nove cursos ( 52,9 \% do total regional ), dos quais nenhum ficou sem conceito ( SC ) ou com conceitos 3, 4 ou 5.
O conceito modal para as Instituições Privadas na região foi 2, atribuído a cinco cursos ( 55,6 \% ), seguido do conceito 1 que foi atribuído a quatro cursos.
As Instituições Públicas participaram com oito cursos na região Norte ( 47,1 \% do total regional ), metade dos quais obteve conceito 3, o valor modal, seguido por dois cursos ( 33,3 \% ) que receberam conceito 4.
Os dois cursos restantes foram alocados nos conceitos 2 e 5 ( um curso, cada ).
Nenhum curso ficou sem conceito ( SC ) ou recebeu o conceito 1 nesta combinação de categoria e região.
Na região Nordeste, a rede privada concentrou 20 dos 47 cursos participantes, o equivalente a 42,6 \% do total da região, a menor das proporções desta categoria numa região.
Dentre os 20 cursos oferecidos por IES Privadas no Nordeste, 15 receberam conceito 2, conceito modal.
Dois cursos ficaram alocados ao conceito 1, outros dois, ao conceito 3 e um único curso, ao conceito 4.
Nenhum curso ficou sem conceito ( SC ) ou obteve conceito 5.
As Instituições Públicas no Nordeste participaram com 27 cursos ( 57,4 \% do total da região ), a maior das proporções desta categoria numa região.
Desses 27 cursos, 14 ( 51,9 \% ) obtiveram conceito 4, o valor modal, seguido por nove cursos que receberam conceito 3 ( 33,3 \% ).
Dois cursos foram alocados ao conceito 2 e outros dois, ao conceito 5.
Nenhum curso ficou sem conceito ( SC ) ou recebeu conceito 1 nesta combinação de categoria e região.
Na região Sudeste, como já mencionado no Capítulo 2, a proporção de cursos da rede privada, 80,3 \%, foi a mais elevada dentre as regiões brasileiras, correspondendo a 114 dos 142 cursos participantes.
Para os cursos oferecidos em IES Privadas, na região Sudeste, o conceito modal foi 2, atribuído a 55 cursos ( 48,2 \% ).
Foi nesta combinação de Categoria Administrativa e Grande Região que dois cursos não receberam conceito ( SC ).
Nenhum curso recebeu conceito 5 e os demais foram assim avaliados : conceito 1 ( 14 cursos ), conceito 3 ( 36 cursos ) e conceito 4 ( sete cursos ).
Entre os 28 cursos oferecidos por Instituições Públicas na região Sudeste, o conceito modal foi 4 ( 13 cursos ), e os demais cursos receberam os conceitos 2 ( dois cursos ), 3 ( nove cursos ) e 5 ( quatro cursos ).
Na região Sudeste, nesta categoria, nenhum curso ficou sem conceito ( SC ) ou recebeu conceito 1.
As Instituições Privadas concentraram 35 dos 55 cursos participantes da região Sul, 63,6 \% do total regional.
Desses, 16 obtiveram conceito 3, o conceito modal, seguidos por outros 15 avaliados com conceito 2.
Nesta combinação de Categoria Administrativa e Grande Região, apenas um curso recebeu conceito 5, e nenhum curso ficou sem conceito ( SC ) ou recebeu conceito 1.
As Instituições Públicas na região Sul participaram com 20 cursos ( 36,4 \% ).
O conceito modal foi 4, atribuído a 11 cursos.
Dois cursos receberam conceito 2, e outros dois, conceito 5.
Os demais cursos receberam conceito 1 ( um curso ) e 3 ( quatro cursos ).
Nenhum curso vinculado a IES Pública ficou sem conceito ( SC ) nessa região.
Na região Centro-Oeste, 14 dos 22 cursos participantes eram de Instituições Privadas ( 63,6 \% em termos regionais ).
Para esses 14 cursos, os conceitos 2 e 3 foram os conceitos modais, tendo sido atribuídos, cada qual, a seis cursos.
Dois cursos receberam conceito 1, e nenhum curso de IES Privada ficou sem conceito ( SC ) ou recebeu os conceitos 4 e 5 nessa região.
Dos oito cursos oferecidos por Instituições Públicas, quatro foram avaliados com conceito 4, conceito modal, dois cursos obtiveram conceito 5 e outros dois foram avaliados com conceito 3.
Nesta região, nenhum curso de IES Pública ficou sem conceito ( SC ) ou recebeu os conceitos 1 ou 2.
O Gráfico 5.2 apresenta a distribuição do Conceito Enade segundo a Categoria Administrativa da IES.
Os cursos em IES Públicas ( linha azul ) apresentam uma poligonal mais à direita do que os em IES Privadas ( linha verde ) e, consequentemente, uma distribuição de Conceitos Enade com valores maiores. 5.3 CONCEITOS POR ORGANIZAÇÃO ACADÊMICA E POR GRANDE REGIÃO Na Tabela 5.3, encontra -se a distribuição dos conceitos atribuídos aos cursos participantes do ENADE/2014 na Área de Engenharia Civil, por Organização Acadêmica, segundo as Grandes Regiões brasileiras.
Como constatado no Capítulo 2, a maioria dos cursos de Engenharia Civil participantes do ENADE/2014 eram oferecidos por Universidades.
Dos 283 cursos aqui analisados, 169 eram oferecidos em Universidades, 55 em Centros Universitários e 59 em Faculdades.
Esta distribuição corresponde a, respectivamente, 59,7 \%, 19,4 \% e 20,8 \% dos cursos.
De acordo com os dados apresentados, de todos os 12 cursos avaliados com conceito 5, nove eram vinculados a Universidades.
Um dos dois cursos sem conceito ( SC ) na Área era oferecido em Universidade.
Esse tipo de Organização Acadêmica teve o conceito 3 como modal, atribuído a 56 cursos ( 33,1 \% ).
Os demais cursos de Universidades avaliados receberam os conceitos 1 ( nove cursos ), 2 ( 45 cursos ), 4 ( 49 cursos ) e conceito 5 ( nove cursos, como já mencionado ).
Entre os cursos oferecidos por Centros Universitários, o conceito modal foi 2, atribuído a 26 cursos ( 47,3 \% ).
Neste tipo de Organização Acadêmica, nenhum curso ficou sem conceito ( SC ), e os outros cursos receberam os conceitos 1 ( seis cursos ), 3 ( 18 cursos ), 4 ( quatro cursos ) e 5 ( um curso ).
Dentre os 59 cursos oferecidos Faculdades avaliados um ficou sem conceito ( SC ). 32 cursos ( 54,2 \% ) obtiveram conceito 2, conceito modal, seguidos de 14 cursos que receberam conceito 3 e de oito cursos aos quais foi atribuído conceito 1.
Os demais cursos, neste tipo de Organização Acadêmica, receberam os conceitos 4 e 5 ( dois cursos cada ).
Considerando -se separadamente as regiões brasileiras, verifica -se que, na região Norte, as Universidades concentraram dez dos 17 cursos participantes.
Quatro cursos ficaram alocados ao conceito 3, conceito modal, seguidos de dois cursos que obtiveram conceito 2 e de outros dois que obtiveram conceito 4.
Um curso recebeu conceito 1 e outro, conceito 5.
Nenhum curso ficou sem conceito ( SC ) nesta categoria e região.
Os Centros Universitários da região Norte foram representados por quatro cursos, três dos quais receberam conceito 2, conceito modal.
O outro curso recebeu conceito 1.
As Faculdades participaram com três cursos na região Norte, os quais receberam conceitos 1 ( dois cursos, conceito modal ) e 2 ( um curso ).
Na região Nordeste, as Universidades participaram com 34 dos 47 cursos na Área de Engenharia Civil.
O conceito modal foi 4, atribuído a 14 cursos, enquanto os conceitos 3 e 2 foram atribuídos, respectivamente, a nove e a oito cursos.
Um curso recebeu conceito 1, e dois cursos receberam conceito 5.
Nenhum dos cursos oferecidos por Universidades no Nordeste ficou sem conceito ( SC ).
Os Centros Universitários contaram com cinco cursos participantes na região Nordeste.
Desses, três cursos receberam conceito 2.
Os dois cursos restantes obtiveram conceitos 3 e 4 ( um curso cada ).
As Faculdades foram representadas por oito cursos na região Nordeste, seis dos quais receberam conceito 2, conceito modal.
Entre os dois cursos restantes, um recebeu conceito 1 e o outro, conceito 3.
Na região Sudeste, as Universidades concentraram 69 dos 142 cursos da região.
Entre os cursos em Universidades, nessa região, o conceito modal foi 3, atribuído a 24 cursos.
Os demais cursos receberam os conceitos 1 ( seis cursos ), 2 ( 21 cursos ), 4 ( 15 cursos ) e 5 ( dois cursos ).
Como já comentado, um único curso oferecido em Universidade ficou sem conceito ( SC ) na região Sudeste.
Os Centros Universitários participaram com 33 cursos na região Sudeste, dos quais 13 obtiveram o conceito modal, 2, seguidos de outros 12 cursos que receberam conceito 3.
Cinco cursos obtiveram conceito 1, e três cursos, conceito 4.
Nenhum curso ficou sem conceito ( SC ) ou recebeu conceito 5.
As Faculdades foram representadas por 40 cursos na região Sudeste, 23 dos quais receberam conceito 2, conceito modal.
Os demais cursos se distribuíram nos conceitos 1 ( três cursos ), 3 ( nove cursos ), 4 ( dois cursos ), 5 ( dois cursos ) e um curso ficou sem conceito ( SC ).
Dos 55 cursos da região Sul, 43 eram de Universidades, para os quais o conceito modal foi 3, atribuído a 15 cursos, seguidos de outros 14 cursos que receberam conceito 4.
Os demais cursos receberam os conceitos 1 ( um curso ), 2 ( 11 cursos ) e 5 ( dois cursos ).
Nesse tipo de organização, nenhum dos cursos ficou sem conceito ( SC ).
Os Centros Universitários da região Sul foram representados por oito cursos, cinco dos quais obtiveram o conceito 2, o conceito modal.
Aos três cursos restantes foram atribuídos os conceitos 3 ( dois cursos ) e 5 ( um curso ).
Na região Sul, para os quatro cursos vinculados a Faculdades o resultado foi : três receberam conceito 2, conceito modal, e um recebeu conceito 1.
Nenhum ficou sem conceito ( SC ) ou recebeu os demais conceitos nessa categoria e região.
Na região Centro-Oeste, 13 dos 22 cursos eram vinculados a Universidades.
Nesse tipo de organização, nenhum curso ficou sem conceito ( SC ) ou com conceito 1.
Os conceitos modais foram os conceitos 3 e 4, atribuídos, cada qual, a quatro cursos.
Três cursos receberam conceito 2, e outros dois, conceito 5.
Os Centros Universitários da região Centro-Oeste contaram com cinco cursos, que obtiveram conceitos 2 ( dois cursos ) e 3 ( três cursos ).
Nenhum curso ficou sem conceito ( SC ) ou recebeu os conceitos restantes na categoria.
Dos quatro cursos oferecidos em Faculdades na região Centro-Oeste, metade ficou alocado no conceito 1, conceito modal.
A outra metade recebeu os conceitos 2 e 3 ( um curso cada ).
O Gráfico 5.3 apresenta a distribuição do Conceito Enade segundo a Organização Acadêmica da IES.
Os cursos em Universidades ( linha azul ) apresentam uma poligonal mais à direita do que as demais, denotando uma distribuição com Conceitos Enade melhores.
No outro extremo, temos os cursos oferecidos em Faculdades ( linha vermelha ) com a poligonal mais à esquerda e, consequentemente, uma distribuição de Conceitos Enade com valores menores.
Cursos de Centros Universitários ( linha verde ) apresentam uma situação intermediária.

% matched lemmas: observação
\end{textsample}
